%% ma: L12-B14-0042
%% lop: 12
%% chuong: 5
%% bai: 14
%% dang: 12.14.19
%% loai: TN4
%% muc: TH
%% dap_an: "D"
%% nguon: "CĐ2.5.pdf giáo viên gửi 27/10/2026 (M07, Lê Minh Khôi – Nguyễn Anh Khôi, Chương II Chủ đề 5 Tích có hướng của hai vectơ), Câu 26"
%% kien_thuc: "Bài 14 (tích có hướng, phương trình mặt phẳng, khoảng cách từ điểm đến mặt phẳng), lớp 11 (thể tích khối chóp, khối hộp)"
%% trang_thai: cho-duyet
%% ly_do: "Dạng mới 12.14.19 đề xuất, chờ giáo viên duyệt"
\begin{ex}
Trong không gian $Oxyz$, cho ba vectơ $\vec a=(-1;1;0)$, $\vec b=(1;1;0)$, $\vec c=(1;1;1)$. Cho hình hộp $OABC.O'A'B'C'$ thỏa mãn $\overrightarrow{OA}=\vec a$, $\overrightarrow{OB}=\vec b$, $\overrightarrow{OC'}=\vec c$. Thể tích của hình hộp $OABC.O'A'B'C'$ bằng
\choice{$\dfrac13$}{$\dfrac23$}{$4$}{\True $2$}
\loigiai{Đáy $OABC$ là hình bình hành nên $\overrightarrow{OC}=\overrightarrow{OB}-\overrightarrow{OA}=(2;0;0)$, $\overrightarrow{OO'}=\overrightarrow{OC'}-\overrightarrow{OC}=(-1;1;1)$.

$\left[\overrightarrow{OA},\overrightarrow{OC}\right]=(0;0;-2)$: $S_{OABC}=2$, mặt đáy $(OABC)$: $z=0$, chiều cao $d\left(O';(OABC)\right)=1$. Vậy $V=2$.}
\end{ex}
